%%%%%%%%%%%%%%%%%%%%%%%%%%%%%%%%%%%%%%%%%%%%%%%%%%%%%%%%%%%%%%%%%%%%%%
% Title:    bibupdate manual
% Author:   Andrew Mathas
% Created:  November 2026
% Modified: 15:54 Monday, 02 November 2026
%%%%%%%%%%%%%%%%%%%%%%%%%%%%%%%%%%%%%%%%%%%%%%%%%%%%%%%%%%%%%%%%%%%%%%
\documentclass{amsart}

%%%%%%%%%%%%%%%%%%%%%%%%%%%%%%%%%%%%%%%%%%%%%%%%%%%%%%%%%%%%%%%%%%%%%%
\usepackage[
  layout       = a4paper,
  footskip     = 16mm,
  hmargin      = 16mm,
  includefoot,
  top          = 0.05\paperheight,
  bottom       = 0.12\paperheight,
]{geometry}

\usepackage[svgnames]{xcolor}
\usepackage{booktabs}
\usepackage{fvextra}

% BEGIN GENERATED BIBUPDATE PREAMBLE
\usepackage[T1]{fontenc}
\providecommand{\tightlist}{\setlength{\itemsep}{0pt}\setlength{\parskip}{0pt}}
\makeatletter
\newcommand{\BibUpdateOptionHeading}{\@startsection{paragraph}{4}{0pt}
  {-6pt plus -2pt minus -1pt}{3pt}
  {\normalfont\bfseries\color{NavyBlue}}}
\makeatother
\newenvironment{BibUpdateOptionDescription}
  {\begin{list}{}{\setlength{\leftmargin}{2em}\setlength{\rightmargin}{0pt}%
    \setlength{\listparindent}{0pt}%
    \setlength{\topsep}{0pt}\setlength{\parsep}{4pt plus 1pt}%
    \setlength{\itemsep}{0pt}}\item\relax}
  {\end{list}}
% END GENERATED BIBUPDATE PREAMBLE

\usepackage{imakeidx}
\indexsetup{level=\section*, toclevel=section, noclearpage}
\makeindex[intoc,columns=3]

\newcommand\option[2]{%
  \medskip\noindent\textsf{\textcolor{NavyBlue}{#1}} (default: #2)%
  \index{#1}%
}

\newcommand\ctan[1]{\href{https://www.ctan.org/pkg/#1}{\texttt{#1}}}

\ExplSyntaxOn
\prop_new:N \l__bibupdate_prop

\cs_new_protected:Npn \__bibupdate_read_line:w #1 = #2 \q_stop
{
  \prop_put:Nee \l__bibupdate_prop { \tl_trim_spaces:n {#1} } { \tl_trim_spaces:n {#2} }
}

\ior_new:N \l__bibupdate_ior
\ior_open:Nn \l__bibupdate_ior { bibupdate.ini }
\ior_str_map_inline:Nn \l__bibupdate_ior
{
  \tl_if_in:nnT {#1} { = } { \__bibupdate_read_line:w #1 \q_stop }
}
\ior_close:N \l__bibupdate_ior

\NewExpandableDocumentCommand\BibUpdateIni{m}{ \prop_item:Nn \l__bibupdate_prop {#1} }
\ExplSyntaxOff


%%%%%%%%%%%%%%%%%%%%%%%%%%%%%%%%%%%%%%%%%%%%%%%%%%%%%%%%%%%%%%%%%%%%%%
\newcommand\Section[1]{\subsection*{\textcolor{blue}{#1}}}

\makeatletter
% README-style headings: left aligned, unnumbered, and on their own line.
\renewcommand\subsection{\@startsection{subsection}{2}{0pt}
  {-12pt plus -3pt minus -2pt}{6pt}
  {\normalfont\large\bfseries\color{NavyBlue}}}
\renewcommand\subsubsection{\@startsection{subsubsection}{3}{0pt}
  {-8pt plus -2pt minus -1pt}{3pt}
  {\normalfont\bfseries\color{NavyBlue}}}
\author{Andrew Mathas}
\usepackage{tikz}
\usetikzlibrary{shadows.blur}
\tikzset{
    shadowed/.style={
        blur shadow={shadow blur steps=5},
        bottom color=Teal!80,
        draw=NavyBlue!70,shade,
      % font=\normalfont\Huge\bfseries\scshape,
        rounded corners=8pt,
        top color=MidnightBlue,
      },
      boxes/.style={
        draw=SteelBlue,
        fill=LightSkyBlue!20,
        font=\sffamily\small,
        inner sep=5pt,
        rectangle,
        rounded corners=8pt,
        text=RoyalBlue,
     }
}

\usepackage{hologo}
\usepackage{xspace}
\newcommand\BibUpdate{\textsc{Bib}\textsf{Update}\xspace}
\newcommand\BibUpdateTitle
{%
  \begin{tikzpicture}[remember picture,overlay]
      \node[yshift=-3cm] at (current page.north west)
        {\begin{tikzpicture}[remember picture, overlay]
          \draw[shadowed](30mm,0) rectangle node[white]{\BibUpdate} (\paperwidth-30mm,16mm);
          \node[anchor=west,boxes] at (4cm,0cm) {Andrew Mathas};
          \node[anchor=east,boxes] at (\paperwidth-4cm,0) {Version \BibUpdateIni{version}};
         \end{tikzpicture}
        };
   \end{tikzpicture}
   \vspace*{20mm}
}

\def\@evenfoot{\thepage\hfill\BibUpdate}
\def\@oddfoot {Version \BibUpdateIni{version} --- \BibUpdateIni{release date}\hfill\thepage}
\def\@evenhead{}
\def\@oddhead{}

\usepackage[colorlinks=true,linkcolor=MidnightBlue,urlcolor=SlateBlue]{hyperref}
\hypersetup{
    pdfcreator = { Generated by pdfLaTeX },
    pdfinfo    = {
      Author   = { \BibUpdateIni{author} },
      Keywords = { \BibUpdateIni{keywords} },
      License  = { \BibUpdateIni{license} },
      Subject  = { \BibUpdateIni{description} },
      Title    = { bibupdate - \BibUpdateIni{version} }
    },
}
\makeatother


\begin{document}
\raggedbottom
\setlength{\parindent}{0pt}
\setlength{\parskip}{4pt plus 1pt}

    \BibUpdateTitle


% BEGIN GENERATED BIBUPDATE HELP
% Generated from bibupdate.__doc__; edit the help, then run python3 setup.py readme.
Automatically update the entries of a bibtex file

Usage:

\begin{Verbatim}[breaklines=true,fontsize=\small]
bibupdate [-h|-H] [-a] [-X] [-c] [-d] [-k] [-i FIELDS] [-l LOG]
          [-m | -M | -n] [-o] [-q] [--threads N] [-w LEN]
          [--duplicates | --duplicate-papers | --duplicate-keys]
          [--list-duplicates | --list-papers | --list-keys]
          bibtexfile [outputfile]
\end{Verbatim}

This is a command line tool for updating the entries in a
\href{http://www.bibtex.org/}{BibTeX} file using the free tools provided
by the American Mathematical Society:
\href{https://mathscinet.ams.org/mathscinet/freetools/mrlookup}{mrlookup},
\href{https://mathscinet.ams.org/mathscinet/freetools/mrlookup}{MathSciNet}
and \href{https://mathscinet.ams.org/mathscinet/freetools/mref}{MRef}.
By default \href{https://github.com/AndrewMathas/BibUpdate}{bibupdate}
uses
\href{https://mathscinet.ams.org/mathscinet/freetools/mrlookup}{mrlookup}
to update \texttt{article}, \texttt{book}, \texttt{inproceedings} and
\texttt{incollection} entries that have no \texttt{mrnumber} field. An
\texttt{mrnumber} field, even an empty one, skips this lookup unless
\texttt{-\/-all} or \texttt{-\/-check} is used. In addition, with the
\texttt{-X} option,
\href{https://github.com/AndrewMathas/BibUpdate}{bibupdate} can query
the \href{http://arxiv.org/}{arXiv} and add \texttt{eprint} and
\texttt{archiveprefix} fields to any entry that it finds there.

\textbf{Options by task}:

\textbf{Updating and checking}

\begin{Verbatim}[breaklines=true,fontsize=\small]
-m, --mrlookup        use mrlookup to update bibtex entries (default)
-M, --mathscinet      use mathscinet to update bibtex entries (less flexible)
-n, --mref            use mref to update bibtex entries
-a, --all             update or validate ALL BibTeX entries
-X, --arxiv           also add arXiv eprint fields to entries
-c, --check           check all entries and report duplicates
--threads N          number of concurrent AMS lookups (default: 1)
\end{Verbatim}

\textbf{Duplicate listing}

\begin{Verbatim}[breaklines=true,fontsize=\small]
-d, --list-duplicates list duplicate papers and keys without writing a file
--list-papers        list duplicate papers only
--list-keys          list repeated citation keys only
\end{Verbatim}

\textbf{Duplicate cleanup}

\begin{Verbatim}[breaklines=true,fontsize=\small]
--duplicates         remove duplicate papers and rename repeated keys
--duplicate-papers   remove duplicate papers, keeping the first
--duplicate-keys     rename repeated keys using authors and year
\end{Verbatim}

\textbf{Output and formatting}

\begin{Verbatim}[breaklines=true,fontsize=\small]
-o, --overwrite       overwrite existing bibtex file
-i FIELDS, --ignored-fields FIELDS
                      a string of bibtex fields to ignore
-k, --keep_fonts      do NOT replace fonts \Bbb, \germ and \scr in titles
-w LEN, --wrap LEN    wrap bibtex fields to specified width
\end{Verbatim}

\textbf{Reports and logging}

\begin{Verbatim}[breaklines=true,fontsize=\small]
-q, --quiet           brief updates and missed lookups (repeat to suppress)
-l LOG, --log LOG     log messages to specified file (defaults to stdout)
\end{Verbatim}

\textbf{Help and version}

\begin{Verbatim}[breaklines=true,fontsize=\small]
-h, --help            show grouped options and exit
-H, --Help            print full program description
-v, --version         print the program version and exit
\end{Verbatim}

\textbf{Note:} As described below, you should check the new file for
errors before deleting the original version of your
\href{http://www.bibtex.org/}{BibTeX} file.

By default, \href{https://github.com/AndrewMathas/BibUpdate}{bibupdate}
does not change your original database file. Instead, it creates a new
file with the name \emph{updated\_file.bib}, if your original file was
\emph{file.bib}. It is also possible to have it replace your current
file (use carefully!), or to specify an output file.

\href{http://www.bibtex.org/}{BibTeX} is widely used by the
\href{http://en.wikipedia.org/wiki/LaTeX}{LaTeX} community to maintain
publication databases. This script attempts to add missing fields to the
papers in a \href{http://www.bibtex.org/}{BibTeX} database file by
querying
\href{https://mathscinet.ams.org/mathscinet/freetools/mrlookup}{mrlookup}
and getting the missing information from there. This is not completely
routine because to search on
\href{https://mathscinet.ams.org/mathscinet/freetools/mrlookup}{mrlookup}
you need either the authors or the title of the article and both of
these can have non-standard representations. If the article is already
published then it is also possible to use the publication year and its
page numbers. To search on
\href{https://mathscinet.ams.org/mathscinet/freetools/mrlookup}{mrlookup}
we:

\begin{itemize}
\tightlist
\item
  use the first author\textquotesingle s surname (accents and compound
  surnames can be problematic)
\item
  use numeric page ranges or a single numeric page, when available for a
  published entry
\item
  include the year for books, or when a numeric first page is available
\item
  use the cleaned title when no numeric first page is available
\end{itemize}

The author query uses only the first author\textquotesingle s surname:
\href{https://mathscinet.ams.org/mathscinet/freetools/mrlookup}{mrlookup}
does not interpret a \href{http://www.bibtex.org/}{BibTeX} list of
authors joined by \texttt{and} as separate names. Braces used to protect
capitals in titles (for example \texttt{\{Kac-Moody\}}) are removed for
searching and matching; the title stored in the entry retains its
braces.

If there is a unique match from
\href{https://mathscinet.ams.org/mathscinet/freetools/mrlookup}{mrlookup}
then \href{https://github.com/AndrewMathas/BibUpdate}{bibupdate} updates
the fields supplied by
\href{https://mathscinet.ams.org/mathscinet/freetools/mrlookup}{mrlookup},
except for the citation key and the fields excluded by
\texttt{-\/-ignored-fields}. Local fields that are not supplied by
\href{https://mathscinet.ams.org/mathscinet/freetools/mrlookup}{mrlookup},
such as \texttt{eprint}, are retained. By default, a message is printed
whenever existing fields in the database are changed. If the title of
the retrieved paper does not (fuzzily) match that of the original
article then the entry is NOT updated and a warning message is printed.

Although some care is taken to make sure that the new
\href{http://www.bibtex.org/}{BibTeX} entries correspond to the same
paper that the original entry referred to there is always a (small?)
chance the new entry corresponds to an entirely different paper. In my
experience this happens rarely, and mostly with unpublished manuscripts.
In any case, before you delete your original
\href{http://www.bibtex.org/}{BibTeX} file \emph{you are strongly
advised to check the updated file BibTeX file carefully for errors!}

By default the program prints each updated entry\textquotesingle s
citation key on its own line, followed by an indented bibliography-style
reference. Report lines wrap at 78 columns, preserving the complete
authors, title and publication details. If
\href{http://arxiv.org/}{arXiv} fields are added as well, they appear in
the same reference. \texttt{-\/-check} uses this compact layout to show
entries that would change, without writing an updated file. You can also
compare the original and updated files using \emph{diff}.

As \href{https://github.com/AndrewMathas/BibUpdate}{bibupdate} calls
\href{https://mathscinet.ams.org/mathscinet/freetools/mrlookup}{mrlookup}
this program will only be useful if you have papers in your database
that are listed in
\href{https://mathscinet.ams.org/mathscinet/freetools/mrlookup}{MathSciNet}.
As described below it is also possible to call
\href{https://mathscinet.ams.org/mathscinet/freetools/mrlookup}{MathSciNet}
directly, however, this is less flexible because the \texttt{mrnumber}
field for each paper is required.

I wrote this script because I wanted to automatically add links to
journals, the \href{http://arxiv.org/}{arXiv} and DOIs to the
bibliographies of my papers using
\href{http://www.ctan.org/pkg/hyperref}{hyperref}. This script allowed
me to add the missing urls and DOI fields to my
\href{http://www.bibtex.org/}{BibTeX} database. As a bonus the script
helped me to correct many minor errors that had crept into my
\href{http://www.bibtex.org/}{BibTeX} file over the years (for example,
incorrect page numbers and publication years). Now I use the program to
automatically update the preprint entries in my database when the papers
appear in
\href{https://mathscinet.ams.org/mathscinet/freetools/mrlookup}{MathSciNet}
after they are published.

\subsection*{Options and defaults}\label{options-and-defaults}

Without options,
\href{https://github.com/AndrewMathas/BibUpdate}{bibupdate} uses
\href{https://mathscinet.ams.org/mathscinet/freetools/mrlookup}{mrlookup}
with one AMS lookup at a time, prints compact references to standard
output, and writes \texttt{updated\_file.bib} beside the input
\texttt{file.bib}. The \texttt{-\/-all}, \texttt{-\/-arxiv},
\texttt{-\/-check}, \texttt{-\/-duplicates}, \texttt{-\/-keep\_fonts}
and \texttt{-\/-overwrite} switches are off. Field-value wrapping is off
(\texttt{-\/-wrap\ 0}).

\subsubsection*{Updating and checking}\label{updating-and-checking}

\BibUpdateOptionHeading*{\texorpdfstring{\texttt{-m,\ -\/-mrlookup}}{-m, -\/-mrlookup}}\label{-m---mrlookup}

\begin{BibUpdateOptionDescription}

Use mrlookup to update bibtex entries (default)

\end{BibUpdateOptionDescription}

\BibUpdateOptionHeading*{\texorpdfstring{\texttt{-M,\ -\/-mathscinet}}{-M, -\/-mathscinet}}\label{-m---mathscinet}

\begin{BibUpdateOptionDescription}

Use mathscinet to update bibtex entries

\end{BibUpdateOptionDescription}

\BibUpdateOptionHeading*{\texorpdfstring{\texttt{-n,\ -\/-mref}}{-n, -\/-mref}}\label{-n---mref}

\begin{BibUpdateOptionDescription}

Use mref to update bibtex entries

These three options select which of the free AMS tools is used to update
the \href{http://www.bibtex.org/}{BibTeX} entries in the database. They
are all free services provided by the American Mathematical Society and
\href{https://github.com/AndrewMathas/BibUpdate}{bibupdate} talks to
them through the JSON web APIs behind the AMS
\href{https://mathscinet.ams.org/mathscinet/freetools}{free tools}
pages.

By default
\href{https://mathscinet.ams.org/mathscinet/freetools/mrlookup}{mrlookup}
is used. Its advantage is that it allows flexible searching by author,
title, year and page numbers, so it can find papers even when they do
not yet have an \texttt{mrnumber}.

With \texttt{-M} the entry is looked up directly by its
\texttt{mrnumber}. A field containing digits is required: missing, empty
or non-numeric MR numbers are skipped, including with \texttt{-\/-all}
and \texttt{-\/-check}.

With \texttt{-n} the whole entry is handed to
\href{https://mathscinet.ams.org/mathscinet/freetools/mref}{MRef} as a
free-form reference. Its best matching record is still checked against
the local title by
\href{https://github.com/AndrewMathas/BibUpdate}{bibupdate}. Like
\href{https://mathscinet.ams.org/mathscinet/freetools/mrlookup}{mrlookup},
\href{https://mathscinet.ams.org/mathscinet/freetools/mref}{MRef}
normally checks only \texttt{article}, \texttt{book},
\texttt{inproceedings} and \texttt{incollection} entries without an
\texttt{mrnumber} field; use \texttt{-\/-all} to check all recognised
entry types.

\end{BibUpdateOptionDescription}

\BibUpdateOptionHeading*{\texorpdfstring{\texttt{-a,\ -\/-all}}{-a, -\/-all}}\label{-a---all}

\begin{BibUpdateOptionDescription}

Update or validate ALL BibTeX entries

By default this switch is off. With
\href{https://mathscinet.ams.org/mathscinet/freetools/mrlookup}{mrlookup}
or \href{https://mathscinet.ams.org/mathscinet/freetools/mref}{MRef},
enabling it checks all recognised entry types, including entries with an
\texttt{mrnumber} field. With \texttt{-\/-arxiv}, it also checks entries
that already have an \texttt{eprint} field. The \texttt{-M} tool always
requires a usable MR number, even with \texttt{-\/-all}.

\end{BibUpdateOptionDescription}

\BibUpdateOptionHeading*{\texorpdfstring{\texttt{-X,\ -\/-arxiv}}{-X, -\/-arxiv}}\label{-x---arxiv}

\begin{BibUpdateOptionDescription}

Also add arXiv eprint fields to entries

This switch is off by default. In addition to the usual AMS lookup,
query the \href{http://arxiv.org/}{arXiv} for each entry and, when the
paper is found there, add \texttt{eprint} and \texttt{archiveprefix}
fields to the \href{http://www.bibtex.org/}{BibTeX} entry:

\begin{Verbatim}[breaklines=true,fontsize=\small]
eprint = {2211.07069},
archiveprefix = {arXiv}
\end{Verbatim}

\href{https://github.com/AndrewMathas/BibUpdate}{bibupdate} searches the
\href{http://arxiv.org/}{arXiv} using the significant words in the title
and only adds these fields when there is a unique paper whose title is a
good (fuzzy) match. This option works together with any of the
\texttt{-m}, \texttt{-M} or \texttt{-n} lookup tools. The
\href{http://arxiv.org/}{arXiv} lookup changes only \texttt{eprint} and
\texttt{archiveprefix}. Entries with an \texttt{eprint} field, even an
empty one, are skipped unless \texttt{-a} or \texttt{-\/-check} is used.
When such entries are checked, these two fields can be replaced. Entries
without a usable title cannot be searched.

\href{https://github.com/AndrewMathas/BibUpdate}{bibupdate} waits at
least 3 seconds after each completed \href{http://arxiv.org/}{arXiv}
request before starting the next one. Updating a large database with
\texttt{-X} can therefore take a while.

\end{BibUpdateOptionDescription}

\BibUpdateOptionHeading*{\texorpdfstring{\texttt{-c,\ -\/-check,\ -\/-check1}}{-c, -\/-check, -\/-check1}}\label{-c---check---check1}

\begin{BibUpdateOptionDescription}

Check all entries and report duplicates

This switch is off by default. It prints each entry that would change as
a compact reference, using the proposed database fields and retaining
any other local fields. The original
\href{http://www.bibtex.org/}{BibTeX} file is not changed and no output
bibliography is written. \texttt{-\/-check1} is an alias for
\texttt{-\/-check}. Check mode implies \texttt{-\/-all}; it runs
\href{http://arxiv.org/}{arXiv} lookups only when \texttt{-\/-arxiv} is
also supplied. Proposed references remain visible even with
\texttt{-qq}.

Before network lookups, checks the input file for repeated citation keys
and papers listed under different keys. Shared DOI, MR number or arXiv
identifiers identify duplicate papers. Matching normalized titles and
authors are reported as possible duplicates; capitalization braces,
punctuation, spacing and given names versus initials are normalized.
Warnings name the keys and source line numbers for review, and remain
visible with \texttt{-\/-quiet}. Duplicate warnings appear in purple on
terminals; log files and redirected output remain plain text.

\end{BibUpdateOptionDescription}

\BibUpdateOptionHeading*{\texorpdfstring{\texttt{-\/-threads\ N}}{-\/-threads N}}\label{--threads-n}

\begin{BibUpdateOptionDescription}

Number of concurrent AMS lookups

Use \texttt{-\/-threads\ 2} to overlap two AMS lookups. The default is
one thread; the thread count must be a positive integer. This works with
\href{https://mathscinet.ams.org/mathscinet/freetools/mrlookup}{mrlookup},
\href{https://mathscinet.ams.org/mathscinet/freetools/mrlookup}{MathSciNet}
and \href{https://mathscinet.ams.org/mathscinet/freetools/mref}{MRef}.
Entries and log messages retain their original order.

\href{http://arxiv.org/}{arXiv} lookups remain sequential and keep their
delay between requests, regardless of the thread count.

\end{BibUpdateOptionDescription}

\subsubsection*{Duplicate listing}\label{duplicate-listing}

\BibUpdateOptionHeading*{\texorpdfstring{\texttt{-d,\ -\/-list-duplicates}}{-d, -\/-list-duplicates}}\label{-d---list-duplicates}

\begin{BibUpdateOptionDescription}

List repeated papers and citation keys without changing the bibliography

No network queries, output bibliography or backups are made, even with
\texttt{-\/-overwrite} or an output filename. Warnings show keys and
source lines in purple on terminals, remain visible with
\texttt{-\/-quiet}, and are suppressed by \texttt{-qq}. If no duplicates
are found, nothing is printed.

\end{BibUpdateOptionDescription}

\BibUpdateOptionHeading*{\texorpdfstring{\texttt{-\/-list-papers}}{-\/-list-papers}}\label{--list-papers}

\begin{BibUpdateOptionDescription}

List duplicate papers only, including possible title-and-author matches

\end{BibUpdateOptionDescription}

\BibUpdateOptionHeading*{\texorpdfstring{\texttt{-\/-list-keys}}{-\/-list-keys}}\label{--list-keys}

\begin{BibUpdateOptionDescription}

List repeated citation keys only

The same listing-only behaviour and quiet levels apply to all three
listing options. Listing options may be combined to show both kinds of
duplicates.

\end{BibUpdateOptionDescription}

\subsubsection*{Duplicate cleanup}\label{duplicate-cleanup}

\BibUpdateOptionHeading*{\texorpdfstring{\texttt{-\/-duplicates} (alias
\texttt{-\/-deduplicate})}{-\/-duplicates (alias -\/-deduplicate)}}\label{--duplicates-alias---deduplicate}

\begin{BibUpdateOptionDescription}

Remove duplicate papers and rename repeated citation keys

All duplicate modes are off by default. This offline cleanup first
removes repeated papers, then renames repeated keys among the remaining
entries. It makes no AMS or \href{http://arxiv.org/}{arXiv} queries. The
output is \texttt{updated\_file.bib} beside \texttt{file.bib}, or the
supplied output filename. Use \texttt{-\/-overwrite} to replace the
input; existing non-empty output files are backed up to \texttt{.bak} as
usual.

The retained entries keep their original fields and formatting:
ignored-field, font-conversion and wrapping options do not apply to
cleanup. Cleanup cannot be combined with listing options or
\texttt{-\/-check}. Other lookup options are ignored.

\textbf{Changed meaning:} \texttt{-\/-duplicates} now writes a cleaned
bibliography. Use \texttt{-d} or \texttt{-\/-list-duplicates} for the
previous listing-only behaviour.

\end{BibUpdateOptionDescription}

\BibUpdateOptionHeading*{\texorpdfstring{\texttt{-\/-duplicate-papers} (alias
\texttt{-\/-remove-duplicate-papers})}{-\/-duplicate-papers (alias -\/-remove-duplicate-papers)}}\label{--duplicate-papers-alias---remove-duplicate-papers}

\begin{BibUpdateOptionDescription}

Remove repeated papers, keeping the first occurrence in file order

Papers match through a shared DOI, MR number or arXiv identifier, or
through matching normalized titles and authors, as described under
\texttt{-\/-check}. Connected matches form one group, even when linked
through different identifiers. Papers with the same citation key are
also compared. Title-and-author matching is heuristic; inspect
\texttt{-\/-list-papers} before removing entries.

Only the first entry\textquotesingle s fields are kept; fields from
later entries are not merged. Each removal reports the removed key and
the retained key, with source lines. Literal \texttt{crossref},
\texttt{xref}, \texttt{related}, \texttt{entryset} and \texttt{xdata}
references inside the bibliography are updated where the removed key has
no surviving occurrence. String macros are preserved without expansion.
Citations in external LaTeX files are not changed; use the reported
mappings to update citations to removed keys.

\end{BibUpdateOptionDescription}

\BibUpdateOptionHeading*{\texorpdfstring{\texttt{-\/-duplicate-keys} (alias
\texttt{-\/-rename-duplicate-keys})}{-\/-duplicate-keys (alias -\/-rename-duplicate-keys)}}\label{--duplicate-keys-alias---rename-duplicate-keys}

\begin{BibUpdateOptionDescription}

Rename repeated keys without deleting papers

Keep the first occurrence of each key. Later occurrences receive a key
formed from all authors\textquotesingle{} surnames and the year, for
example \texttt{AcharRiche2019}. Editors are used if authors are
missing, and \texttt{Anonymous} if neither is available. When the key is
already used, add \texttt{a}, \texttt{b}, ..., \texttt{z}, \texttt{aa},
etc., reserving all original keys and newly generated keys. Each rename
is reported with its source line; use \texttt{-\/-log} to save the
mappings. References to the original repeated key continue to refer to
its first surviving occurrence; ambiguous citations need manual review.

Using \texttt{-\/-duplicate-papers} together with
\texttt{-\/-duplicate-keys} is equivalent to \texttt{-\/-duplicates}.

\end{BibUpdateOptionDescription}

\subsubsection*{Output and formatting}\label{output-and-formatting}

\BibUpdateOptionHeading*{\texorpdfstring{\texttt{-o,\ -\/-overwrite}}{-o, -\/-overwrite}}\label{-o---overwrite}

\begin{BibUpdateOptionDescription}

Overwrite the existing bibtex file with the updated version

This switch is off by default. Without it, the output is
\texttt{updated\_file.bib} beside the input \texttt{file.bib}, or the
output filename supplied as the last argument. With \texttt{-o}, the
input file is replaced and any output filename is ignored. An existing,
non-empty output file is copied to its \texttt{.bak} file before
writing, including when \texttt{-o} replaces the input. An existing
\texttt{.bak} is replaced.

\end{BibUpdateOptionDescription}

\BibUpdateOptionHeading*{\texorpdfstring{\texttt{-i\ FIELDS,\ -\/-ignored-fields=FIELDS}}{-i FIELDS, -\/-ignored-fields=FIELDS}}\label{-i-fields---ignored-fieldsfields}

\begin{BibUpdateOptionDescription}

A string of \href{http://www.bibtex.org/}{BibTeX} fields to ignore when
writing the updated file

By default \href{https://github.com/AndrewMathas/BibUpdate}{bibupdate}
removes the following fields from each
\href{http://www.bibtex.org/}{BibTeX} entry:

\begin{itemize}
\tightlist
\item
  \texttt{coden}
\item
  \texttt{mrreviewer}
\end{itemize}

The \texttt{-i} option replaces this default list. Repeated \texttt{-i}
options are combined; field names are case-insensitive and separated by
spaces. An empty string keeps every field:

\begin{Verbatim}[breaklines=true,fontsize=\small]
bibupdate -i "coden fjournal" file.bib   # ignore coden and fjournal
bibupdate -i coden -i fjournal file.bib  # ignore coden and fjournal
bibupdate -i "" file.bib                 # do not ignore any fields
\end{Verbatim}

\end{BibUpdateOptionDescription}

\BibUpdateOptionHeading*{\texorpdfstring{\texttt{-k,\ -\/-keep\_fonts}}{-k, -\/-keep\_fonts}}\label{-k---keep_fonts}

\begin{BibUpdateOptionDescription}

do NOT replace fonts \texttt{\textbackslash{}Bbb},
\texttt{\textbackslash{}germ} and \texttt{\textbackslash{}scr} in titles

This switch is off by default, so supported font commands in titles are
replaced. The \href{http://www.bibtex.org/}{BibTeX} entries generated by
\href{https://mathscinet.ams.org/mathscinet/freetools/mrlookup}{mrlookup}
use \texttt{\textbackslash{}Bbb}, \texttt{\textbackslash{}germ} and
\texttt{\textbackslash{}scr} for the \texttt{\textbackslash{}mathbb},
\texttt{\textbackslash{}mathfrak} and \texttt{\textbackslash{}mathscr}
fonts. By default, in the \emph{title} fields, these fonts
specifications are automatically changed to the following more
\href{http://en.wikipedia.org/wiki/LaTeX}{LaTeX} friendly fonts:

\begin{itemize}
\tightlist
\item
  \texttt{\textbackslash{}Bbb\ X} -\/-\textgreater{}
  \texttt{\textbackslash{}mathbb\{X\}}
\item
  \texttt{\textbackslash{}scr\ X} -\/-\textgreater{}
  \texttt{\textbackslash{}mathcal\{X\}}
\item
  \texttt{\textbackslash{}germ\ X} -\/-\textgreater{}
  \texttt{\textbackslash{}mathfrak\{X\}}
\end{itemize}

Use \texttt{-k} to retain the original font commands. This applies both
to input titles and to titles retrieved from AMS.

\end{BibUpdateOptionDescription}

\BibUpdateOptionHeading*{\texorpdfstring{\texttt{-w\ LEN,\ -\/-wrap\ LEN}}{-w LEN, -\/-wrap LEN}}\label{-w-len---wrap-len}

\begin{BibUpdateOptionDescription}

Wrap bibtex fields to specified width

The default is \texttt{0}, which leaves field values unwrapped. A
positive \texttt{LEN} wraps each field value to that width, with a tab
before continuation lines. Field names, braces and indentation add to
the total line length, so \texttt{LEN} is not a strict limit on complete
output lines. Negative values are rejected.

\end{BibUpdateOptionDescription}

\subsubsection*{Reports and logging}\label{reports-and-logging}

\BibUpdateOptionHeading*{\texorpdfstring{\texttt{-q,\ -\/-quiet,\ -\/-quieter}}{-q, -\/-quiet, -\/-quieter}}\label{-q---quiet---quieter}

\begin{BibUpdateOptionDescription}

Print fewer messages

By default, \href{https://github.com/AndrewMathas/BibUpdate}{bibupdate}
prints a wrapped reference for each updated entry and a notice for each
attempted AMS lookup that has no unique matching paper, including
preprints. With \texttt{-q} or \texttt{-\/-quiet} it prints brief update
notices, missed lookups and network warnings, without separator bars or
field-by-field changes. Missed lookups appear in red when the output is
a terminal; redirected output and log files contain plain text. Reports
wrap at 78 columns. \texttt{-\/-quieter} is retained as an alias. Repeat
the option (\texttt{-qq}) to suppress update notices and lookup
warnings. Fatal errors are still reported.

The \texttt{-\/-check} option prints proposed references in the compact
layout, including when combined with \texttt{-q} or \texttt{-qq}.
Duplicate warnings follow the same warning levels as missed lookups. Use
\texttt{-l} to send messages to a log file.

\end{BibUpdateOptionDescription}

\BibUpdateOptionHeading*{\texorpdfstring{\texttt{-l\ LOG,\ -\/-log\ LOG}}{-l LOG, -\/-log LOG}}\label{-l-log---log-log}

\begin{BibUpdateOptionDescription}

Log output to file (defaults to stdout)

By default reports and warnings go to standard output. Use \texttt{-l}
to write them to a file instead; an existing log file is replaced. Fatal
errors are also printed to standard error.

\end{BibUpdateOptionDescription}

\subsubsection*{Help and version}\label{help-and-version}

\BibUpdateOptionHeading*{\texorpdfstring{\texttt{-h,\ -\/-help}}{-h, -\/-help}}\label{-h---help}

\begin{BibUpdateOptionDescription}

Show a short list of options, grouped by task, and exit.

\end{BibUpdateOptionDescription}

\BibUpdateOptionHeading*{\texorpdfstring{\texttt{-H,\ -\/-Help}}{-H, -\/-Help}}\label{-h---help-1}

\begin{BibUpdateOptionDescription}

Print this full description, including grouped options, defaults and
known issues, and exit.

\end{BibUpdateOptionDescription}

\BibUpdateOptionHeading*{\texorpdfstring{\texttt{-v,\ -\/-version}}{-v, -\/-version}}\label{-v---version}

\begin{BibUpdateOptionDescription}

Print the program version and licence, and exit.

\end{BibUpdateOptionDescription}

\subsection*{Known issues}\label{known-issues}

\begin{itemize}
\tightlist
\item
  \textbf{BibTeX syntax:} the parser uses regular expressions rather
  than a full BibTeX grammar. Use brace-delimited entries with one field
  per line and braced values or simple bare words. Multiple fields on
  one line, quoted values, \texttt{\#} concatenation, string-macro
  expansion, and fields containing \texttt{@} or \texttt{=\{...\}} can
  be misread or omitted. Inspect the output carefully when using these
  forms. Duplicate listing and cleanup use a source scanner that
  preserves quoted and nested values, comments and macros; macros are
  not expanded for duplicate matching.
\item
  \textbf{Title matching:} capitalization braces and apostrophes are
  removed before searching and comparison. Dollar-delimited mathematics
  is also removed, which can make distinct titles look alike. Other TeX
  commands and accents are not fully decoded, so equivalent spellings
  can still fail to match. A unique fuzzy title match is a heuristic,
  not a guarantee that the paper is correct.
\item
  \textbf{Author names:}
  \href{https://mathscinet.ams.org/mathscinet/freetools/mrlookup}{mrlookup}
  uses only the first author\textquotesingle s surname. A compound
  surname is preserved in \texttt{surname,\ given\ names} form; without
  a comma, only the last word is used. Accents, name particles and
  alternative spellings can therefore affect searches.
\item
  \textbf{Page numbers:} numeric ranges and single numeric pages are
  supported. Article identifiers and prefixed page numbers are not
  handled reliably. If no numeric first page is extracted, the title is
  used instead; for non-book entries the year is then omitted from the
  query.
\item
  \textbf{Ambiguous or missing results:} an AMS entry is updated only
  when one candidate survives title matching. No match or multiple
  matches leave it unchanged and produce a missed-lookup notice. Entries
  without a usable title have less protection against an incorrect
  match. Skipped entries, such as those with an existing
  \texttt{mrnumber} under the default settings, are not failed lookups
  and produce no missed-lookup notice.
\item
  \textbf{arXiv searches:} searches use significant title words and
  examine at most ten results. Titles dominated by mathematics or short
  words can be missed. No unique title match leaves the arXiv fields
  unchanged; this is normally silent. With \texttt{-\/-all} or
  \texttt{-\/-check}, an existing identifier can be replaced by the
  identifier of the matching title, so review such changes.
\item
  \textbf{Network services:} connections can fail, time out or return
  unexpected data. Requests use a 30-second timeout. Connection errors
  and timeouts produce warnings, and processing continues. A missed
  lookup does not prove that the paper is absent from the database.
\end{itemize}
% END GENERATED BIBUPDATE HELP

\subsection*{Rebuilding the documentation}

  The built-in help in \texttt{bibupdate.py} is the source for \texttt{README.md}
  and the marked, generated regions of this manual. The title page, styling and
  text outside those regions can be edited directly.

  Run \texttt{python3 setup.py readme} in the source directory to refresh both
  sources and compile \texttt{bibupdate.pdf}. This requires Python, Pandoc, and
  \texttt{latexmk} from a TeX distribution. Use \texttt{--no-pdf} to refresh
  only the sources; \texttt{docs} is an alias for \texttt{readme}.
  \BibUpdate itself remains a self-contained Python script with no
  non-standard dependencies.

\subsection*{Feature requests and bug reports}

  Please make any feature requests and bug reports on the \BibUpdate github
  page:
  \begin{center}
    \href{\BibUpdateIni{repository}/issues}{\texttt{\BibUpdateIni{repository}/issues}}.
  \end{center}
  Please provide as much detail as possible when submitting requests.

  Bug reports must be accompanied by a \textit{minimal working
  example}, which is a minimal \hologo{BibTeX} entry that demonstrates the
  problem.

% BEGIN GENERATED BIBUPDATE LICENCE
\subsection*{Licence}\label{licence}

Andrew Mathas

\href{https://github.com/AndrewMathas/BibUpdate}{bibupdate} Version
1.4.0-r0. Copyright (C) 2014-2026

GNU General Public License, Version 3, 29 June 2007

This program is free software: you can redistribute it and/or modify it
under the terms of the GNU General Public License
(\href{http://www.gnu.org/licenses/gpl.html}{GPL}) as published by the
Free Software Foundation, either version 3 of the License, or (at your
option) any later version.

This program is distributed in the hope that it will be useful, but
WITHOUT ANY WARRANTY; without even the implied warranty of
MERCHANTABILITY or FITNESS FOR A PARTICULAR PURPOSE. See the GNU General
Public License for more details.
% END GENERATED BIBUPDATE LICENCE

\end{document}
